% 作者题证的编辑转录副本，不是独立下载的上游原件。
% 来源：https://terrytao.wordpress.com/2011/09/27/254a-notes-3-haar-measure-and-the-peter-weyl-theorem/

\begin{theorem}[Finite-dimensional separation; Tao's Theorem 7]
Let $G$ be a compact Hausdorff group with Haar probability measure $\mu$.
For every $y\ne e$ there is a finite-dimensional left-invariant subspace
$V\subseteq L^2(G,\mu)$ on which the left translation
$\tau(y)f(x)=f(y^{-1}x)$ is not the identity. In particular, continuous
finite-dimensional unitary representations separate the elements of $G$.
\end{theorem}
\paragraph{Convolution operators.}
The regular representation $\tau:G\to U(L^2(G))$ is strongly continuous.
For $g\in L^2(G)$ let
\[
 T_gf=f*g,\qquad (f*g)(x)=\int_G f(z)g(z^{-1}x)\,d\mu(z).
\]
The convolution belongs to $C(G)$. If
\begin{equation}\label{selfadjoint-kernel}
 g(x^{-1})=\overline{g(x)},
\end{equation}
then $T_g$ is self-adjoint. It commutes with all left translations.
Its integral kernel $K(z,x)=g(z^{-1}x)$ belongs to $L^2(G\times G)$;
indeed $\|K\|_2=\|g\|_2$ by invariance and $\mu(G)=1$.
Thus $T_g$ is compact by the compactness theorem for square-integrable
integral kernels. Each nonzero eigenspace is left-invariant.

\begin{lemma}[Compact self-adjoint spectral theorem; Tao's Theorem 6]
A compact self-adjoint operator $T$ on a complex Hilbert space $H$ has
an orthogonal decomposition
\[
 H=V_0\mathbin{\widehat\oplus}\mathop{\widehat\bigoplus}_{n}V_{\lambda_n},
\]
where $V_0=\ker T$, all $\lambda_n$ are nonzero real eigenvalues, and their
eigenspaces are finite-dimensional. The distinct nonzero eigenvalues are
at most countable and, if infinite in number, tend to zero.
\end{lemma}
\begin{proof}
Self-adjointness makes eigenspaces for distinct eigenvalues orthogonal and
forces eigenvalues to be real. For $r>0$, choose an orthonormal basis of
$\bigoplus_{|\lambda|>r}V_\lambda$. The images under $T$ of distinct basis
vectors are orthogonal and have norms at least $r$. This basis cannot be
infinite, since its image would have no convergent subsequence, contrary
to compactness. Consequently the sum is finite-dimensional for every
$r>0$. This proves finite-dimensionality of every nonzero eigenspace and
the asserted countability and convergence of the eigenvalues.

Let $W$ be the orthogonal complement of the closed span of all eigenspaces,
including the kernel. It is invariant under $T$, and $T|_W$ is self-adjoint
and has no eigenvectors. Suppose $W\ne0$. Then $\|T|_W\|>0$ because the
restriction has trivial kernel. For a bounded self-adjoint operator,
\[
 \|T|_W\|=\sup_{x\in W,\,\|x\|\leq1}|\langle Tx,x\rangle|.
\]
Choose $x_n\in W$, $\|x_n\|\leq1$, such that
$\langle Tx_n,x_n\rangle\to\lambda$, where
$\lambda=\|T|_W\|$ or $\lambda=-\|T|_W\|$. Then
\[
\begin{split}
 0\leq\|Tx_n-\lambda x_n\|^2
 &=\|Tx_n\|^2+\lambda^2\|x_n\|^2
       -2\lambda\langle Tx_n,x_n\rangle\\
 &\leq 2\lambda^2-2\lambda\langle Tx_n,x_n\rangle\longrightarrow0.
\end{split}
\]
Applying $T$ gives $T(Tx_n)-\lambda Tx_n\to0$. By compactness, pass to a
subsequence with $Tx_n\to v$. Then $Tv=\lambda v$. Since $T|_W$ has no
eigenvectors, $v=0$. But this implies
$\langle Tx_n,x_n\rangle\to0$, a contradiction to $\lambda\ne0$.
Thus $W=0$.
\end{proof}

\begin{proof}[Proof of finite-dimensional separation]
Suppose, for a contradiction, that $\tau(y)$ is the identity on every
finite-dimensional invariant subspace. For each $g$ satisfying
\eqref{selfadjoint-kernel}, the nonzero eigenspaces of $T_g$ are such
subspaces. The operator $\tau(y)-1$ annihilates them. It also preserves
$\ker T_g$, since it commutes with $T_g$. The spectral decomposition
therefore gives
\[
 T_g(\tau(y)-1)f=0\qquad(f\in L^2(G)).
\]
Equivalently,
\[
 \tau(y)(f*g)=f*g
\]
for every $f,g\in L^2(G)$ with $g$ satisfying \eqref{selfadjoint-kernel}.
Choose an open symmetric identity neighbourhood $U$ so small that
$y\notin U^2$, and put $f=g=\mathbf1_U$. These functions are square
integrable, $g$ satisfies \eqref{selfadjoint-kernel}, and
\[
 (f*g)(e)=\mu(U)>0,\qquad (f*g)(y)=0.
\]
Since $f*g$ is continuous, the asserted translation equality holds
everywhere, and at $y$ would imply $(f*g)(e)=(f*g)(y)$, a contradiction.
Restriction of the regular representation to the resulting
finite-dimensional invariant subspace is a continuous unitary
representation with nontrivial value at $y$.
\end{proof}
